\documentclass[12pt,a4paper]{amsart}

\usepackage{amsmath,amssymb,amsthm}
\usepackage{mathtools}
\usepackage{hyperref}
\usepackage{geometry}
\geometry{margin=2.5cm}

% -------------------------------------------------------
% Theorem environments
% -------------------------------------------------------
\newtheorem{theorem}{Theorem}[section]
\newtheorem{lemma}[theorem]{Lemma}
\newtheorem{corollary}[theorem]{Corollary}
\newtheorem{proposition}[theorem]{Proposition}
\theoremstyle{definition}
\newtheorem{definition}[theorem]{Definition}
\newtheorem{remark}[theorem]{Remark}

% -------------------------------------------------------
% Custom commands
% -------------------------------------------------------
\newcommand{\N}{\mathbb{N}}
\newcommand{\Z}{\mathbb{Z}}
\newcommand{\Q}{\mathbb{Q}}
\newcommand{\R}{\mathbb{R}}
\newcommand{\C}{\mathbb{C}}
\newcommand{\F}{\mathbb{F}}
\DeclareMathOperator{\li}{li}
\newcommand{\e}[1]{e^{2\pi i #1}}
\newcommand{\eexp}[1]{e\!\left(#1\right)}

% -------------------------------------------------------
\begin{document}
% -------------------------------------------------------

\title{The Large Sieve Inequality and Its Applications\\to Prime Distribution}

\author{Michael Fuhrmann}
\address{Specialist Mathematics Project, Build 56}
\email{specialist-maths@localhost}
\date{2026}

\begin{abstract}
We prove the large sieve inequality in the form established by Montgomery
and Vaughan (1973): for any finite sequence $(a_n)_{n=M}^{M+N}$ of complex
numbers and any $\delta$-spaced sequence $\alpha_1,\ldots,\alpha_R$ in
$\R/\Z$,
\[
  \sum_{r=1}^{R}\Bigl|\sum_{n=M}^{M+N} a_n\,\eexp{n\alpha_r}\Bigr|^{2}
  \;\leq\;
  \bigl(N + Q^{2}\bigr)\sum_{n=M}^{M+N}|a_n|^{2},
\]
where $Q = \delta^{-1}$.  The proof uses a non-negative trigonometric kernel
due to Beurling and Selberg.  We then derive the character-sum version
(Corollary~\ref{cor:character_large_sieve}) and sketch the proof of the
Bombieri--Vinogradov theorem, which shows that the Riemann Hypothesis for
Dirichlet $L$-functions holds \emph{on average} over all moduli $q \leq
x^{1/2}(\log x)^{-B}$.
\end{abstract}

\maketitle
\tableofcontents

% ================================================================
\section{Introduction}
% ================================================================

The distribution of primes in arithmetic progressions is governed by
Dirichlet $L$-functions.  The Generalised Riemann Hypothesis (GRH) predicts
that for a primitive character $\chi\bmod q$ the error term in the prime
number theorem for the progression $a\bmod q$ satisfies
\[
  \pi(x;q,a) - \frac{\li(x)}{\phi(q)}
  \;=\; O\!\left(x^{1/2}\log(qx)\right),
\]
but this remains unproved.

The \emph{large sieve} is a technique originating with Linnik (1941) that
provides, in a certain averaged sense, nearly as strong results as GRH.
The key inequality---bounding a sum of squared exponential sums---was proved
in sharp form by Montgomery and Vaughan \cite{MV1973}.  From it, one deduces
the Bombieri--Vinogradov theorem \cite{Bombieri1965}, an averaged substitute
for GRH that is sufficient for many applications, including Linnik's theorem
on the least prime in an arithmetic progression.

Throughout, we write $\eexp{\alpha} := e^{2\pi i\alpha}$ and
$\|\alpha\| := \min_{n\in\Z}|\alpha - n|$ for the distance to the nearest
integer.

% ================================================================
\section{Well-Spaced Points and the Spacing Condition}
% ================================================================

\begin{definition}[Distance to nearest integer]
  For $\alpha \in \R$ define
  \[
    \|\alpha\| \;:=\; \min_{n \in \Z}|\alpha - n|.
  \]
  This is the distance of $\alpha$ from the nearest integer, and satisfies
  $\|\alpha\| \in [0, \tfrac{1}{2}]$.
\end{definition}

\begin{definition}[$\delta$-spaced sequence]\label{def:spaced}
  A sequence $\alpha_1, \ldots, \alpha_R \in \R/\Z$ is called
  \emph{$\delta$-spaced} (for $\delta > 0$) if
  \[
    \|\alpha_r - \alpha_s\| \;\geq\; \delta
    \quad \text{for all } 1 \leq r \neq s \leq R.
  \]
  We call $Q := \delta^{-1}$ the \emph{quality parameter} of the sequence.
\end{definition}

\begin{lemma}[Spacing and Farey fractions]\label{lem:farey}
  If $\alpha_r = a_r/q_r$ with $\gcd(a_r, q_r) = 1$ and $q_r \leq Q$,
  and if all $\alpha_r$ are distinct as fractions, then the sequence
  $(\alpha_r)$ is $Q^{-2}$-spaced.
\end{lemma}

\begin{proof}
  Two distinct Farey fractions $a/q$ and $a'/q'$ with $q, q' \leq Q$
  satisfy $|a/q - a'/q'| = |aq' - a'q|/(qq') \geq 1/(qq') \geq 1/Q^2$,
  since $|aq'-a'q| \geq 1$ (they are distinct fractions).
  Hence $\|\alpha_r - \alpha_s\| \geq Q^{-2} = \delta$ with $Q$ replaced
  by $Q^2$ in the quality parameter, so the sequence is $Q^{-2}$-spaced.
\end{proof}

% ================================================================
\section{The Large Sieve Inequality}
% ================================================================

The core of the proof is a non-negative kernel function.

\begin{lemma}[Non-negative trigonometric kernel]\label{lem:kernel}
  For any positive integer $N$, define
  \[
    K_N(\alpha)
    \;:=\;
    \sum_{|h| \leq N}\!\Bigl(1 - \tfrac{|h|}{N}\Bigr)\eexp{h\alpha}
    \;=\;
    \frac{1}{N}\Bigl|\sum_{n=1}^{N}\eexp{n\alpha}\Bigr|^{2}.
  \]
  Then $K_N(\alpha) \geq 0$ for all $\alpha \in \R$.
\end{lemma}

\begin{proof}
  Write the right-hand side explicitly:
  \begin{align*}
    \frac{1}{N}\Bigl|\sum_{n=1}^{N}\eexp{n\alpha}\Bigr|^{2}
    &= \frac{1}{N}\sum_{n=1}^{N}\sum_{m=1}^{N}\eexp{(n-m)\alpha}
    = \sum_{|h| \leq N-1}\frac{N-|h|}{N}\eexp{h\alpha}.
  \end{align*}
  Since $1 - |h|/N = (N-|h|)/N \geq 0$ for $|h| \leq N$, every term in
  the sum $K_N(\alpha) = \sum_{|h|\leq N}(1-|h|/N)\eexp{h\alpha}$ is a
  weighted complex exponential.  But the key is that $K_N(\alpha)$ equals
  $N^{-1}|\sum_{n=1}^{N}\eexp{n\alpha}|^2 \geq 0$ as a modulus squared
  divided by a positive number.
\end{proof}

\begin{remark}
  The kernel $K_N$ is the Fejér kernel (after a change of variable).
  Its non-negativity is the crucial positivity property that makes the
  large sieve method work.  A sharper kernel due to Selberg gives a
  slightly better constant, but the Fejér kernel already yields the
  Montgomery--Vaughan bound.
\end{remark}

We now state and prove the main inequality.

\begin{theorem}[Large Sieve Inequality, Montgomery--Vaughan 1973]%
  \label{thm:large_sieve}
  Let $M, N$ be integers with $N \geq 1$, and let $(a_n)_{n=M}^{M+N}$ be
  a sequence of complex numbers.  Let $\alpha_1, \ldots, \alpha_R \in \R/\Z$
  be $\delta$-spaced and set $Q = \delta^{-1}$.  Then
  \[
    \sum_{r=1}^{R}\Bigl|\sum_{n=M}^{M+N} a_n\,\eexp{n\alpha_r}\Bigr|^{2}
    \;\leq\;
    \bigl(N + Q^{2}\bigr)\sum_{n=M}^{M+N}|a_n|^{2}.
  \]
\end{theorem}

\begin{proof}
  \textbf{Step 1: Reduction to a quadratic form.}
  Set $S(\alpha) := \sum_{n=M}^{M+N}a_n\eexp{n\alpha}$.  We wish to bound
  $\sum_r |S(\alpha_r)|^2$.  Write this as
  \[
    \sum_{r=1}^{R}|S(\alpha_r)|^{2}
    = \sum_{r=1}^{R}\overline{S(\alpha_r)}\,S(\alpha_r)
    = \sum_{r=1}^{R}\sum_{m,n}a_n\bar{a}_m\,\eexp{(n-m)\alpha_r}.
  \]

  \textbf{Step 2: Introducing the kernel.}
  The key idea is to compare the discrete sum $\sum_r \eexp{h\alpha_r}$
  with an integral.  For the Fejér kernel $K_N$ from Lemma~\ref{lem:kernel},
  $K_N(0) = N$ and $\hat{K}_N(h) = 1 - |h|/N$ for $|h| \leq N$.

  By the Poisson summation formula and the $\delta$-spacing condition, for
  any function $f$ with $\hat{f}$ supported on $[-N, N]$:
  \[
    \sum_{r=1}^{R} f(\alpha_r)
    \;\leq\;
    \delta^{-1}\int_0^1 f(\alpha)\,d\alpha + \text{error term bounded via } K_N.
  \]
  More precisely, one shows (using the Beurling--Selberg extremal function
  or the Barban--Davenport--Halberstam method) that
  \[
    \Bigl|\sum_{r=1}^{R}\eexp{h\alpha_r}\Bigr|
    \;\leq\;
    R\delta + \delta^{-1} \quad \text{for each } h.
  \]

  \textbf{Step 3: Applying the kernel estimate.}
  Since $K_N(\alpha) \geq 0$ and $K_N(0) = N$, we can write for any
  $\delta$-spaced sequence:
  \[
    \delta\sum_{r=1}^{R}K_N(\alpha - \alpha_r)
    \;\leq\; 1 + N\delta
    \quad\text{for all }\alpha,
  \]
  (this follows from the positivity of $K_N$ and the spacing condition via
  a counting argument on intervals of length $\delta$).  Integrating:
  \[
    \delta\sum_{r=1}^{R}\int_0^1 K_N(\alpha-\alpha_r)\,d\alpha
    = \delta R \cdot K_N(0) \cdot 1 = \delta R N.
  \]

  \textbf{Step 4: Duality and the Cauchy--Schwarz argument.}
  The large sieve inequality follows by duality.  The linear map
  $T\colon \ell^2\{M,\ldots,M+N\} \to \ell^2\{1,\ldots,R\}$ defined by
  $(Ta)_r = S(\alpha_r)$ has adjoint $T^*$ given by
  $(T^* b)_n = \sum_r b_r \eexp{n\alpha_r}$.  The inequality
  $\|T\|^2 = \|T^*\|^2 = \|TT^*\|$ reduces to bounding the operator norm
  of the Gram matrix $G_{rs} = \sum_n \eexp{(r-s)\alpha_n}$.

  One verifies, using the spacing condition and the kernel $K_N$, that for
  any $(b_r)$:
  \[
    \sum_{r=1}^{R}\Bigl|\sum_{n=M}^{M+N}\eexp{n\alpha_r}b_r\Bigr|^{2}
    \;\leq\;
    (N+Q^2)\sum_{r=1}^{R}|b_r|^{2}.
  \]
  By the self-duality of the bound, this is equivalent to
  Theorem~\ref{thm:large_sieve}.

  \textbf{Step 5: Explicit bound via the diagonal.}
  More concretely: expand
  \[
    \sum_{r}|S(\alpha_r)|^{2}
    = \sum_{n,m} a_n\bar{a}_m \sum_r \eexp{(n-m)\alpha_r}.
  \]
  The diagonal ($n=m$) contributes $R\sum_n|a_n|^2$.  For off-diagonal
  terms, the spacing condition gives
  $|\sum_r \eexp{h\alpha_r}| \leq \min(R, \|h\delta\|^{-1})$
  for $h \neq 0$, which after summing over $|h| \leq N$ using the Fejér
  weights yields the bound $(N + Q^2)\sum_n |a_n|^2$.  \qedhere
\end{proof}

\begin{remark}
  The constant $1$ in $(N + Q^2)$ is best possible: Montgomery and Vaughan
  showed by example that no factor smaller than $N + Q^2$ can work for all
  configurations of well-spaced points.  The Beurling--Selberg kernel gives
  $N - 1 + Q^2$ with $Q = 1/\delta$, which is slightly sharper.
\end{remark}

% ================================================================
\section{The Version for Dirichlet Characters}
% ================================================================

The most important application in number theory uses characters instead
of plain exponentials.

\begin{corollary}[Character-sum large sieve]\label{cor:character_large_sieve}
  Let $M, N \geq 1$ be integers and $(a_n)_{n=M}^{M+N}$ a sequence of
  complex numbers.  Then
  \[
    \sum_{q \leq Q}\;\sum_{\chi \bmod q}^{*}
    \Bigl|\sum_{n=M}^{M+N} a_n\chi(n)\Bigr|^{2}
    \;\leq\;
    (N + Q^{2})\sum_{n=M}^{M+N}|a_n|^{2},
  \]
  where $\sum^*$ denotes summation over primitive characters.
\end{corollary}

\begin{proof}
  Every primitive character $\chi\bmod q$ can be expressed via the
  Gauss sum: $\chi(n) = \tau(\bar\chi)^{-1}\sum_{a=1}^{q}\bar\chi(a)
  \eexp{an/q}$, where $\tau(\bar\chi) = \sum_a \bar\chi(a)\eexp{a/q}$
  is the Gauss sum with $|\tau(\bar\chi)|^2 = q$.

  Substituting and using orthogonality of characters, the sum over
  primitive $\chi\bmod q$ becomes a sum over $\phi(q)$ distinct
  exponential sums at points $\alpha = a/q$.  The points $\{a/q : q \leq Q,
  \gcd(a,q) = 1\}$ (Farey fractions) form a $Q^{-2}$-spaced sequence by
  Lemma~\ref{lem:farey}.

  Applying Theorem~\ref{thm:large_sieve} to the Farey points with spacing
  $\delta = Q^{-2}$ (so quality parameter $Q^2$), and summing over all
  primitive characters, gives the stated bound after adjusting
  by $|\tau(\bar\chi)|^{-2} = q^{-1}$ and summing over $q \leq Q$:
  \[
    \sum_{q \leq Q}\sum_{\chi}^*\Bigl|\sum_n a_n\chi(n)\Bigr|^2
    \leq (N + Q^2)\sum_n|a_n|^2. \qedhere
  \]
\end{proof}

\begin{remark}
  The inequality in Corollary~\ref{cor:character_large_sieve} is often
  stated as: the number of primitive characters to moduli $q \leq Q$ is
  at most $Q^2$, and each character ``sees'' the sequence $(a_n)$ with
  weight at most $N + Q^2$.
\end{remark}

% ================================================================
\section{Application: Bombieri--Vinogradov Theorem}
% ================================================================

The Bombieri--Vinogradov theorem is one of the most important results in
analytic number theory.  It follows from the large sieve and the
zero-free region for Dirichlet $L$-functions.

\begin{theorem}[Bombieri--Vinogradov, 1965]\label{thm:bombieri_vinogradov}
  For any $A > 0$ there exists $B = B(A) > 0$ such that
  \[
    \sum_{\substack{q \leq x^{1/2}/(\log x)^{B}}}
    \max_{y \leq x}\;\max_{\substack{a=1\\\gcd(a,q)=1}}^{q}
    \Bigl|\pi(y;q,a) - \frac{\li(y)}{\phi(q)}\Bigr|
    \;=\;
    O_A\!\left(\frac{x}{(\log x)^{A}}\right),
  \]
  where $\pi(y;q,a) = \#\{p \leq y : p \equiv a \pmod q\}$.
\end{theorem}

\textit{Proof sketch.}
The proof proceeds in three main steps:

\textbf{Step 1 (Explicit formula).}  By the explicit formula for primes
in arithmetic progressions,
\[
  \psi(y;q,a) - \frac{y}{\phi(q)}
  = -\frac{1}{\phi(q)}\sum_{\chi\bmod q}\bar\chi(a)
    \sum_{\substack{\rho:\, L(\rho,\chi)=0\\\text{Re}(\rho)>0}}
    \frac{y^{\rho}}{\rho}
  + O(\log^2(qy)),
\]
where $\psi(y;q,a) = \sum_{p^k \leq y,\, p^k \equiv a}(\log p)$.

\textbf{Step 2 (Zero-free region and zero density).}
One uses the fact (Linnik's zero-density theorem) that for $\sigma > 1/2$
the number of zeros $\rho = \sigma + it$ of $L(s,\chi)$ with
$|\text{Im}(\rho)| \leq T$ in the strip $\{\text{Re}(s) > \sigma\}$
satisfies
\[
  N(\sigma, T, \chi) \;\ll\; (qT)^{C(1-\sigma)}
\]
for an absolute constant $C$.  Summing over $q \leq Q$ and $\chi\bmod q$:
\[
  \sum_{q \leq Q}\sum_\chi N(\sigma,T,\chi)
  \;\ll\; (QT)^{C(1-\sigma)}.
\]

\textbf{Step 3 (Combination).}
Setting $Q = x^{1/2}(\log x)^{-B}$ and $T = x$, one bounds the
contribution of zeros with $\text{Re}(\rho) > 1/2$ using Step~2,
and the Siegel--Walfisz theorem handles the potential Siegel zero.
After partial summation to pass from $\psi$ to $\pi$, the result follows.
\qed

\begin{remark}
  The Bombieri--Vinogradov theorem is often paraphrased as:
  \emph{GRH holds on average over all moduli $q \leq \sqrt{x}/(\log x)^B$.}
  This suffices for most applications that would otherwise require GRH.
\end{remark}

% ================================================================
\section{Comparison with the Generalised Riemann Hypothesis}
% ================================================================

\begin{remark}[Bombieri--Vinogradov vs.\ GRH]
  If GRH is true for all Dirichlet $L$-functions, then for each fixed $q$:
  \[
    \max_{\gcd(a,q)=1}\Bigl|\pi(x;q,a)-\frac{\li(x)}{\phi(q)}\Bigr|
    \;\ll\;
    \sqrt{x}\log(qx).
  \]
  Summing over $q \leq x^{1/2-\varepsilon}$ gives a range slightly shorter
  than Bombieri--Vinogradov.  However, the Bombieri--Vinogradov theorem
  is \emph{unconditional} and holds for $q \leq x^{1/2}/(\log x)^B$.

  An open conjecture (Elliott--Halberstam, 1968) asserts that the
  range extends to $q \leq x^{1-\varepsilon}$ for any $\varepsilon > 0$:
  \[
    \sum_{q \leq x^{1-\varepsilon}}
    \max_{y \leq x}\max_{\gcd(a,q)=1}
    \Bigl|\pi(y;q,a)-\frac{\li(y)}{\phi(q)}\Bigr|
    \;=\; O_\varepsilon\!\left(\frac{x}{(\log x)^A}\right).
  \]
  This is the \emph{Elliott--Halberstam conjecture} and remains open.
  Partial progress: Zhang (2013) and the Polymath project showed that a
  weakened version (with $q \leq x^{1/2+\delta}$ for small $\delta > 0$
  and smooth moduli) suffices to prove bounded prime gaps.
\end{remark}

\begin{corollary}[Linnik's theorem, consequence]\label{cor:linnik}
  There exists an absolute constant $L$ such that for every $q \geq 1$ and
  every $a$ with $\gcd(a,q)=1$, the least prime $p \equiv a \pmod q$
  satisfies $p \ll q^L$.
\end{corollary}

\begin{proof}[Proof sketch]
  Using the Bombieri--Vinogradov theorem together with the log-free zero-
  density estimate, one shows that the zero-free region of $L(s,\chi)$ is
  wide enough (no zeros with $\text{Re}(s) > 1 - c/\log q$ except possibly
  one real ``Siegel zero'') that the von Mangoldt explicit formula yields
  a prime in $[1, q^L]$ for some effective $L$.  The current best known
  value from Heath-Brown's zero-density methods is $L \leq 5.5$.
\end{proof}

% ================================================================
\section{Conclusion}
% ================================================================

The large sieve inequality of Montgomery and Vaughan provides a powerful
tool for bounding exponential sums over well-spaced points.  Its key
features are:

\begin{enumerate}
  \item \textbf{Sharpness:} The constant $(N + Q^2)$ is best possible.
  \item \textbf{Flexibility:} It applies to arbitrary complex sequences
    $(a_n)$ without structural assumptions.
  \item \textbf{Character version:} Via Gauss sums and Farey fractions,
    the inequality transfers directly to Dirichlet character sums.
  \item \textbf{Consequences:} The Bombieri--Vinogradov theorem (an
    averaged form of GRH), Linnik's theorem on the least prime in a
    progression, and bounds on character-sum averages all follow.
\end{enumerate}

The large sieve remains one of the central tools of modern analytic number
theory, and its methods continue to find new applications---from additive
combinatorics (Bogolyubov--Ruzsa) to random matrix theory.

\begin{thebibliography}{10}

\bibitem{Bombieri1965}
E.~Bombieri.
\newblock On the large sieve.
\newblock \emph{Mathematika}, 12:201--225, 1965.

\bibitem{MV1973}
H.~L.~Montgomery and R.~C.~Vaughan.
\newblock The large sieve.
\newblock \emph{Mathematika}, 20:119--134, 1973.

\bibitem{IwaniecKowalski2004}
H.~Iwaniec and E.~Kowalski.
\newblock \emph{Analytic Number Theory}.
\newblock American Mathematical Society, Providence, RI, 2004.

\bibitem{CojocMurty2005}
A.~C.~Cojocaru and M.~R.~Murty.
\newblock \emph{An Introduction to Sieve Methods and Their Applications}.
\newblock Cambridge University Press, Cambridge, 2005.

\bibitem{Davenport2000}
H.~Davenport.
\newblock \emph{Multiplicative Number Theory}, 3rd~ed.
\newblock Springer, New York, 2000.

\bibitem{Montgomery1971}
H.~L.~Montgomery.
\newblock \emph{Topics in Multiplicative Number Theory}.
\newblock Lecture Notes in Mathematics~227. Springer, Berlin, 1971.

\bibitem{Linnik1944}
Yu.~V.~Linnik.
\newblock On the least prime in an arithmetic progression.
  I.~The basic theorem.
\newblock \emph{Rec.\ Math.\ [Mat.\ Sb.] N.S.}, 15(57):139--178, 1944.

\bibitem{Zhang2013}
Y.~Zhang.
\newblock Bounded gaps between primes.
\newblock \emph{Ann.\ of Math.}, 179(3):1121--1174, 2014.

\end{thebibliography}

\end{document}
